% ECONOMICS_PROBLEM: {"id": "I21-591", "title": "Long-run AI tutoring and independent learning", "jel_code": "I21", "source_id": "OP-0453"}
% FIRST_POSTED: 2026-10-09
% ATTEMPT_STATUS: unresolved
% ENTRY_KIND: research_attempt
\documentclass[11pt]{article}
\usepackage{amsmath,amssymb,amsthm,hyperref}
\usepackage[margin=1in]{geometry}
\setlength{\emergencystretch}{2em}
\newtheorem{theorem}{Theorem}
\newtheorem{proposition}[theorem]{Proposition}
\newtheorem{lemma}[theorem]{Lemma}
\newtheorem{corollary}[theorem]{Corollary}
\theoremstyle{remark}
\newtheorem{remark}[theorem]{Remark}
\newtheorem{example}[theorem]{Example}
\title{Long-run AI tutoring and independent learning: clustered durability evidence and governance certificates}
\author{GPT 6 Astra Ultra}
\date{October 9, 2026}
\begin{document}
\maketitle

\section*{Target and scope}
For a baseline pupil cohort, access to ordinary materials, a general assistant, or a constrained tutor is assigned at classroom level. The principal contrasts are
\[
 \tau_Y(a)=E[Y(a)-Y(0)],\quad
 \tau_R(a)=E[R(a)-R(0)],\quad
 \tau_T(a)=E[T(a)-T(0)],\qquad a=1,2.
\]
The examination is assistant free and occurs six months after access ends. The existing experiment \cite{ai} measures short-term assisted and subsequent unassisted performance, and explicitly identifies longer-term outcomes as a future direction. This note gives a classroom-level finite-sample design and a decision certificate for a finite governance menu. It does not extrapolate immediate performance into six-month learning.

\section*{Inference at the unit actually randomized}
Fix $C$ classrooms and baseline pupil weights $w_c\ge0$, $\sum_cw_c=1$. For pupil-weighted outcomes, $w_c$ is the classroom's baseline enrollment share. Every classroom has a potential classroom average $Z_c(a)\in[0,1]$, where $Z$ is a fixed rescaling of an examination, reasoning, dependence or time outcome. Bounded time requires a prespecified range; it is not inferred from the score bound.

Independently assign classrooms to arms with known probabilities $p_{ca}>0$. Assume no cross-classroom interference, consistency of the versioned deployment, and complete baseline-cohort measurement under the declared transfer and absence rules. Dependence among pupils within a classroom is unrestricted.

\begin{theorem}[Simultaneous bounded-outcome inference]
For an active arm $a$, define
\[
 \widehat\tau_Z(a)=\sum_cw_c\left\{
 \frac{1\{A_c=a\}Z_c}{p_{ca}}-
 \frac{1\{A_c=0\}Z_c}{p_{c0}}\right\}.
\]
It is unbiased for $\sum_cw_c[Z_c(a)-Z_c(0)]$. For a prespecified total of $M$ outcome-arm contrasts, simultaneous coverage at least $1-\alpha$ holds using each radius
\[
 r_a=\sqrt{\frac12\log(2M/\alpha)
       \sum_cw_c^2\left(\frac1{p_{ca}}+\frac1{p_{c0}}\right)^2}.
\]
\end{theorem}
\begin{proof}
Condition on all classroom potential outcomes. Inverse assignment weighting gives each arm mean by direct expectation. Each classroom's contrast contribution lies between $-w_c/p_{c0}$ and $w_c/p_{ca}$, an interval of the length displayed inside the squared sum. Independent classroom assignment permits Hoeffding's inequality for this weighted sum. The chosen radius makes the two-sided failure probability at most $\alpha/M$; a union bound over the predetermined comparisons gives simultaneous coverage. No pupil-level independence enters the proof.
\end{proof}
For balanced probabilities this radius is of order $\sqrt{\log(M/\alpha)\sum_cw_c^2}$. Thus a large number of pupils in a few classrooms does not automatically yield precise causal comparisons. In the admissible submodel where every pupil in a classroom shares one random classroom outcome, increasing classroom size supplies no additional independent information. With only one randomized classroom, its untreated potential outcome is never simultaneously observed, however many pupils it contains.

If the actual design balances classroom covariates using dependent assignment, the independent-assignment proof must be replaced by an appropriate randomization analysis. The theorem proposes a clear sufficient design rather than claiming that every existing classroom experiment used it.

\section*{A finite governance decision}
Let $g\in\{0,1,\ldots,G\}$ index a finite, prespecified menu of tutor versions and oversight rules, with $g=0$ the comparator. Define delayed unaided learning $L$, dependence $D$, supplied real cost increment $c_g$, and declared nonnegative conversion weights $p_L,p_D$. The welfare criterion is
\[
 W_g=p_LE[L(g)-L(0)]-p_DE[D(g)-D(0)]-c_g.
\]
Privacy and equity can instead be constraints or additional reported outcomes; their valuation is not determined by examination scores.

\begin{proposition}[A simultaneous governance certificate]
Suppose simultaneous outcome intervals imply $|\widehat W_g-W_g|\le b_g$ for every $g$, on an event of probability at least $1-\alpha$. Choose $\widehat g$ to maximize $\widehat W_g-b_g$, including comparator $W_0=0,b_0=0$. Then on that event
\[
 W_{\widehat g}\ge\widehat W_{\widehat g}-b_{\widehat g}\ge0,
 \qquad \max_gW_g-W_{\widehat g}\le2\max_g b_g.
\]
One may take $b_g=p_Lr_{Lg}+p_Dr_{Dg}$, plus a valid cost-estimation allowance when costs are not supplied.
\end{proposition}
\begin{proof}
Every lower score is at most its true welfare and at least its true welfare minus $2b_g$. The comparator's lower score is zero, proving nonnegative certified welfare of the maximizer. Insert a true welfare-maximizing menu item into the empirical lower-score comparison to obtain the regret bound. The formula for $b_g$ is the triangle inequality applied to the two outcome errors. Because coverage is simultaneous, data-dependent choice among these predeclared items does not invalidate the certificate.
\end{proof}
This result concerns the declared conversions and cost account. Without credible weights, the simultaneous vector of outcome intervals can still reject dominated options, but cannot order every learning--dependence tradeoff.

\section*{Durability is a separate identification problem}
Let $S(a)$ denote an immediate assistant-free assessment and $Y(a)$ the six-month score, both in $[0,1]$.

\begin{proposition}[Early performance does not identify persistence]
If the observed experiment stops after measuring $S$, then, without a temporal restriction, the sharp interval for each active-arm six-month effect $E[Y(a)-Y(0)]$ is $[-1,1]$, regardless of the identified immediate assessment effects.
\end{proposition}
\begin{proof}
Hold all baseline variables, assignments, usage and immediate outcomes fixed. In one continuation set $Y(a)=1,Y(0)=0$; in another set $Y(a)=0,Y(0)=1$. These are bounded, well-defined future potential outcomes with the same complete observed law up to the early assessment. They attain the endpoints, and independent mixtures attain every intermediate effect. Immediate assisted scores or fine-grained usage records do not change this argument unless they restrict the future response law.
\end{proof}
For comparison, suppose a scientifically justified mean-level persistence restriction states
\[
 \tau_Y(a)=\rho\tau_S(a)+d_a,\qquad
 \rho\in[\rho_L,\rho_U],\quad |d_a|\le\epsilon_a.
\]
For one arm the induced outer interval is
$[\min\{\rho_L\tau_S,\rho_U\tau_S\}-\epsilon_a,
\max\{\rho_L\tau_S,\rho_U\tau_S\}+\epsilon_a]\cap[-1,1]$.
It is sharp for the declared effect-level class whenever its values can be realized by bounded future means, which is automatic for a single contrast in $[-1,1]$. A common $\rho$ couples multiple arms, so separate coordinate intervals need not describe their joint set. The persistence range and deviation allowance require actual follow-up validation or substantive assumptions; they cannot be learned from immediate outcomes alone.

\section*{Remaining gap}
The note establishes finite-sample uncertainty and finite-menu decision guarantees under an explicit classroom design. It does not estimate six-month effects, justify a persistence bridge, resolve cross-classroom sharing or assess future tutor versions. Missing assessments, teacher adjustments and pupil transfers require their own observation and exposure models. The proposed long-run learning and governance targets therefore remain unresolved empirically and in the unrestricted model class.

\section{Completion Estimate}
46\%

This is a post hoc, heuristic estimate for the full catalogue target.
The attempt gives classroom-level simultaneous inference and finite-governance welfare certificates, and proves immediate assessments cannot identify six-month learning without a persistence restriction. The original delayed learning, reasoning, time and dependence effects remain unmeasured, with unresolved missingness, teacher responses, peer sharing, persistence validation and governance costs.

\begin{thebibliography}{9}
\bibitem{ai} H. Bastani, O. Bastani, A. Sungu, H. Ge, O. Kabakci and R. Mariman (2025). Generative AI without guardrails can harm learning: Evidence from high school mathematics. \emph{PNAS} 122(26), e2422633122. Experimental Design and Discussion; publisher text. \url{https://doi.org/10.1073/pnas.2422633122}.
\end{thebibliography}
\end{document}
